%% ============================================================
%%  Public Health and AI — Beamer presentation
%%  Albert Jan van Hoek
%%  Compile with: pdflatex (twice) or xelatex on Overleaf
%% ============================================================

\documentclass[aspectratio=169, 10pt]{beamer}

% ── Theme ────────────────────────────────────────────────────
\usetheme{Madrid}
\usecolortheme{default}

% ── Custom colour palette ────────────────────────────────────
%  Deep teal primary, off-white background, warm amber accent
\definecolor{primary}{HTML}{065A82}     % deep teal
\definecolor{accent}{HTML}{F4A261}      % warm amber
\definecolor{lightbg}{HTML}{F5F9FC}     % near-white
\definecolor{darktext}{HTML}{1A1A2E}    % near-black

\setbeamercolor{palette primary}{bg=primary, fg=white}
\setbeamercolor{palette secondary}{bg=primary!80, fg=white}
\setbeamercolor{palette tertiary}{bg=primary!60, fg=white}
\setbeamercolor{palette quaternary}{bg=primary, fg=white}
\setbeamercolor{structure}{fg=primary}
\setbeamercolor{background canvas}{bg=lightbg}
\setbeamercolor{normal text}{fg=darktext}
\setbeamercolor{frametitle}{bg=primary, fg=white}
\setbeamercolor{block title}{bg=primary, fg=white}
\setbeamercolor{block body}{bg=primary!8, fg=darktext}
\setbeamercolor{alerted text}{fg=accent}

% ── Fonts ─────────────────────────────────────────────────────
\usepackage[T1]{fontenc}
\usepackage[utf8]{inputenc}
\usepackage{lmodern}

% ── Useful packages ───────────────────────────────────────────
\usepackage{booktabs}
\usepackage{microtype}
\usepackage{tikz}
\usetikzlibrary{arrows.meta, positioning}

% ── Remove navigation symbols ─────────────────────────────────
\setbeamertemplate{navigation symbols}{}

% ── Footer ────────────────────────────────────────────────────
\setbeamertemplate{footline}{%
  \leavevmode%
  \hbox{%
    \begin{beamercolorbox}[wd=.5\paperwidth, ht=2.5ex, dp=1.5ex, left, leftskip=1em]
      {palette primary}
      \usebeamerfont{author in head/foot}Albert Jan van Hoek — RIVM
    \end{beamercolorbox}%
    \begin{beamercolorbox}[wd=.4\paperwidth, ht=2.5ex, dp=1.5ex, center]
      {palette secondary}
      \usebeamerfont{title in head/foot}\insertshorttitle
    \end{beamercolorbox}%
    \begin{beamercolorbox}[wd=.1\paperwidth, ht=2.5ex, dp=1.5ex, right, rightskip=1em]
      {palette primary}
      \usebeamerfont{date in head/foot}\insertframenumber/\inserttotalframenumber
    \end{beamercolorbox}%
  }%
}

% ── Title metadata ────────────────────────────────────────────
\title[AI \& Public Health]{Artificial Intelligence and Public Health}
\subtitle{Rethinking Infectious Disease Surveillance}
\author{Albert Jan van Hoek}
\institute{RIVM — Dutch National Institute for Public Health and the Environment}
\date{}

% ── Custom command: highlight box ─────────────────────────────
\newcommand{\highlight}[1]{%
  \begin{beamercolorbox}[rounded=true, shadow=true, sep=4pt]{block title}
    \centering\large\textbf{#1}
  \end{beamercolorbox}%
}

% =============================================================
\begin{document}

% ── Title slide ───────────────────────────────────────────────
\begin{frame}[plain]
  \vspace{1.5cm}
  \titlepage
\end{frame}

% ── Who am I ─────────────────────────────────────────────────
\begin{frame}{Who am I?}
  \begin{itemize}
    \item Senior researcher at \textbf{RIVM} (Dutch National Institute for
          Public Health and the Environment)
    \item Principal Investigator, \textbf{Infectieradar} — participatory
          surveillance, part of the European Influenzanet network
    \item Health economist and infectious disease modeller
    \item Working on NLP, agentic AI, and transformer modelling
    \item Field experience in \textit{Gambia, Kenya, and Sierra Leone}
    \item Background: \textbf{Economics / Biology / Philosophy}
  \end{itemize}
\end{frame}

% ── Section: Intelligence ─────────────────────────────────────
\section{What is intelligence?}

\begin{frame}{What is intelligence?}
  \begin{center}
    \large
    \textit{``The faculty of understanding; intellect — a capacity to
    understand.''}\\[0.5em]
    — Oxford Dictionary
  \end{center}

  \bigskip
  \begin{alertblock}{}
    \centering
    But: a capacity that lives \textit{where}, exactly?
  \end{alertblock}
\end{frame}

\begin{frame}{Capacity of what? What is ``Artificial''?}
  \begin{columns}[T]
    \column{0.55\textwidth}
      \begin{block}{A shared structure}
        Both biological and artificial systems are instances of a deeper
        phenomenon:\\[0.5em]
        \textbf{Systems that build expectations, meet reality, and adapt to
        the mismatch.}
      \end{block}

    \column{0.4\textwidth}
      % Placeholder for the image originally in the deck.
      % Replace \fbox{...} with \includegraphics[width=\linewidth]{your_image}
      \begin{center}
        \fbox{\parbox[c][4cm][c]{0.9\linewidth}{%
          \centering\small\color{gray}
          [Insert illustration here]\\
          \texttt{Picture5.jpg}}}
      \end{center}
  \end{columns}
\end{frame}

\begin{frame}{Does this sound familiar?}
  \begin{columns}[T]
    \column{0.48\textwidth}
      \begin{block}{Virus}
        A virus does not predict.  
        But across its variants, it \textit{explores} possibilities — 
        and reality selects what works.
      \end{block}

    \column{0.48\textwidth}
      \begin{block}{My hypothesis}
        Intelligence, learning, and evolution are not separate phenomena.\\[0.4em]
        They are \textbf{the same loop}, expressed at different scales.
      \end{block}
  \end{columns}
\end{frame}

\begin{frame}{Neither is a thing}
  \begin{itemize}
    \item A \textbf{virus} is not a thing.  
          It is a \textit{process} — one that only exists within a network of
          hosts.
    \bigskip
    \item \textbf{Intelligence} is not a thing.  
          It is a \textit{process} — one that only exists within a network of
          interactions.
  \end{itemize}

  \bigskip
  \begin{alertblock}{}
    \centering
    Strip away the metaphor, and the structure is identical.
  \end{alertblock}
\end{frame}

\begin{frame}{No substrate, no process}
  The substrate is not a container.  
  \textit{It determines what is possible.}

  \bigskip
  \begin{tabular}{ll}
    \toprule
    Change this\ldots & \ldots and this changes. \\
    \midrule
    Contact patterns    & The epidemic \\
    Data flow           & AI behaviour \\
    Collaboration patterns & Science \\
    \bottomrule
  \end{tabular}

  \bigskip
  \begin{alertblock}{}
    \centering
    You don't need to change the elements.  
    \textbf{You change the network.}
  \end{alertblock}
\end{frame}

\begin{frame}{The real question}
  \begin{block}{The loop}
    Generate → Test → Update\\[0.3em]
    \ldots needs a substrate to run.  
    \textbf{The network is that substrate.}
  \end{block}

  \bigskip
  \begin{columns}[c]
    \column{0.45\textwidth}
      \begin{center}
        \large\color{gray}
        \textit{``How powerful is the AI?''} {\small(wrong question)}
      \end{center}

    \column{0.5\textwidth}
      \begin{alertblock}{}
        \centering\large
        \textbf{How well does your network learn?}
      \end{alertblock}
  \end{columns}
\end{frame}

\begin{frame}{Intelligence is a verb}
  \begin{enumerate}
    \item Elements interact.
    \item Interactions create patterns.
    \item Reality selects among those patterns.
    \item The system changes.
    \item \textit{And the cycle continues.}
  \end{enumerate}

  \bigskip
  \begin{alertblock}{}
    \centering
    Intelligence is not a state — it is a \textbf{verb}.\\[0.3em]
    New structure is not designed.  
    It \textit{emerges}, is tested, and is reshaped.
  \end{alertblock}
\end{frame}

% ── Section: Three Lessons ────────────────────────────────────
\section{Three lessons}

\begin{frame}{Lesson 1: AI is not a normal tool}
  \begin{block}{The adaptive loop}
    \centering
    \textbf{Generate → Test → Update}\\[0.5em]
    The same loop that drives learning and evolution — now operating at
    \textit{digital speed and scale}.
  \end{block}

  \bigskip
  \begin{alertblock}{}
    \centering
    This is the first tool that does not just \textit{execute}.  
    \textbf{It adapts.}
  \end{alertblock}
\end{frame}

\begin{frame}{Lesson 2: Surveillance must become adaptive}
  \begin{columns}[T]
    \column{0.44\textwidth}
      \begin{block}{Current model}
        Data → Analysis → Decision\\[0.4em]
        A pipeline. Linear. Slow.
      \end{block}

    \column{0.52\textwidth}
      \begin{alertblock}{What we need}
        Observe → Interpret → Act → Update → Repeat\\[0.4em]
        This is not just about AI.  
        It is about \textbf{how we organise ourselves}.
      \end{alertblock}
  \end{columns}
\end{frame}

\begin{frame}{Lesson 3: Collective capability is real}
  \begin{itemize}
    \item A virus cannot transmit by itself.  
          An \textbf{epidemic} is a collective capability — of society and
          pathogen together.
    \item The same is true for an organisation, a surveillance network, a
          public health system.
    \bigskip
    \item Learning as an \textit{individual} builds a \textbf{skillset}.
    \item Learning as a \textit{collective} builds \textbf{capability} — and
          that capability is \textit{non-linear}.
  \end{itemize}

  \bigskip
  \begin{alertblock}{}
    \centering
    We can strengthen it. Or let it degrade.
  \end{alertblock}
\end{frame}

% ── Section: Building at RIVM ─────────────────────────────────
\section{Building a learning network}

\begin{frame}{So what does it look like…}
  \begin{center}
    \Large\textit{…when you actually build a learning network?}
  \end{center}
\end{frame}

\begin{frame}{What I am building at RIVM}
  \begin{columns}[T]
    \column{0.48\textwidth}
      \begin{block}{AI-driven}
        \begin{itemize}
          \item NLP for questionnaire analysis \& data harmonisation
          \item Agentic AI for quality assurance
          \item Benchmarking methodology for AI outputs
          \item Transformer modelling \& synthetic populations
          \item Digital-twin architecture for longitudinal data
        \end{itemize}
      \end{block}

    \column{0.48\textwidth}
      \begin{block}{Human \& collective intelligence}
        \begin{itemize}
          \item FAIR and open data infrastructure
          \item Expanded communication and outreach
          \item Open source software
        \end{itemize}
      \end{block}
  \end{columns}

  \bigskip
  \begin{alertblock}{}
    \centering
    Everything on this list has one purpose: \textbf{make the loop run.}
  \end{alertblock}
\end{frame}

\begin{frame}{Pandemic preparedness as a continuous loop}
  \begin{block}{}
    Preparedness is not a project. It is a \textbf{continuous mission}.
  \end{block}

  \bigskip
  \begin{itemize}
    \item The right frame: a never-ending loop of self-improvement, with
          internal regulation of that improvement.
    \item It is multidimensional — logistics, IT infrastructure, research
          protocols, privacy, communication, contracts, training, and more.
  \end{itemize}

  \bigskip
  \begin{alertblock}{}
    \centering
    \textbf{$k$-cover control:} ensuring the network is always covered,  
    from every direction of failure.
  \end{alertblock}
\end{frame}

\begin{frame}{How to strengthen the network}
  \begin{columns}[T]
    \column{0.44\textwidth}
      \begin{block}{Structural level}
        More interaction.\\
        More collaboration.\\
        More feedback.
      \end{block}

    \column{0.52\textwidth}
      \begin{alertblock}{Personal level — a functional requirement}
        Honesty. Repairing conflict. Patience.\\
        Openness. Transparency. Reciprocity.\\[0.5em]
        \small These are not virtues.  
        They are the conditions under which\\
        \textbf{feedback loops stay intact}.
      \end{alertblock}
  \end{columns}
\end{frame}

% ── Section: What you can do ──────────────────────────────────
\section{What you can do}

\begin{frame}{What you can do}
  \begin{enumerate}
    \item \textbf{Use AI.}  
          Seriously, as much as you can. It extends your adaptive loop.
    \bigskip
    \item \textbf{Think at a higher level.}  
          The work will change; the need for domain expertise will not.
    \bigskip
    \item \textbf{Pair your specialisation} with data literacy, quality
          assurance, and interpretation.  
          Infectious disease knowledge, built over years, cannot be
          replicated quickly.
    \bigskip
    \item \textbf{Collaborate deliberately.}  
          Collective capability is not automatic. It requires maintenance.
  \end{enumerate}
\end{frame}

% ── Conclusion ────────────────────────────────────────────────
\section{Conclusion}

\begin{frame}{Conclusion}
  \begin{itemize}
    \item Intelligence is a \textbf{verb}.
    \item Intelligence is a \textbf{network phenomenon}.
    \item We can optimise this verb — in our systems, our organisations,
          ourselves.
  \end{itemize}

  \bigskip
  \begin{block}{}
    \centering
    Do not underestimate AI. \textbf{Interact with it.}\\[0.4em]
    Think in networks: effects are non-linear, connections matter more than
    nodes, and the loop never stops.
  \end{block}
\end{frame}

% ── Closing slide ─────────────────────────────────────────────
\begin{frame}[plain]
  \begin{center}
    \vspace{1cm}
    {\large\textit{Viral wisdom:}}\\[0.8em]
    {\Large Keep alive what keeps you alive.}\\
    {\large And improve it for the next generation.}\\[1.5em]
    \rule{0.6\textwidth}{0.4pt}\\[1em]
    {\LARGE\textbf{\textit{Disco ergo sum.}}}\\[0.3em]
    {\large I learn, therefore I am.}
  \end{center}
\end{frame}

\end{document}